\documentclass[12pt]{article}
\usepackage[utf8]{inputenc}
\usepackage[T1]{fontenc}
\usepackage{hyperref}
\title{Western Highland Purepecha Language Resources for K-12 Curriculum Development}
\author{Compiled by Claude Sonnet 5.5 (claude-sonnet-5-5) for Andr\'as Kornai}
\date{}
\begin{document}
\maketitle

\section{Introduction}
Western Highland Purepecha (P'urh\'epecha, ISO 639-3: pua; Glottocode west2631) is an indigenous language of Michoac\'an, Mexico, with about 120,000 speakers. Searches found Mexican federal education materials in the language and several scholarly works. The search results did not always make clear which dialect each item represents, and I did not open the PDFs, so their exact titles and grades are not confirmed here.

\section{Primary Educational Resources}
\begin{itemize}
\item SEP/DGEIIB historical editorial collection, document ah\_00097 (a Purepecha-related educational publication from the Direcci\'on General de Educaci\'on Ind\'igena e Intercultural Biling\"ue): \url{https://dgeiib.basica.sep.gob.mx/files/fondo-editorial/acervo-historico/ah_00097.pdf}
\item SEP/DGEIIB historical editorial collection, document ah\_00067: \url{https://dgeiib.basica.sep.gob.mx/files/fondo-editorial/acervo-historico/ah_00067.pdf}
\item SEP/DGEIIB intercultural education collection, document cgeib\_00062: \url{https://dgeiib.basica.sep.gob.mx/files/fondo-editorial/educacion-intercultural/cgeib_00062.pdf}
\item ``Dos lenguas para un poema'' (Nueva Escuela Mexicana collection page): \url{https://nuevaescuelamexicana.sep.gob.mx/contenido/coleccion/dos-lenguas-para-un-poema/}
\end{itemize}
The search summary indicated that the SEP has published Purepecha-language free textbooks for primary grades (including third and fourth grade) and a children's dictionary made by pupils in Ch\'eran and Angahuan; which of the PDFs above corresponds to which was not verified.

\section{Linguistic and Reference Materials}
\begin{itemize}
\item Monz\'on, article on a grammar for P'urhepecha speakers and dialectal variation (journal PDF): \url{https://linguisticamexicana-amla.colmex.mx/index.php/Linguistica_mexicana/article/download/385/355}
\item ``Grammar for P'urhepecha speakers: dialectal variation and standard'' (repository record): \url{https://revistaschilenas.uchile.cl/handle/2250/216755}
\item Redalyc article PDF on Purepecha: \url{https://www.redalyc.org/pdf/1345/134553393006.pdf}
\item Gilberti, \emph{Arte de la lengua de Michoac\'an} (1558, historical grammar, reprint edition): \url{https://linkgua-ediciones.com/en/producto/art-of-the-language-of-michoacan/}
\end{itemize}

\section{Religious and Community Resources}
None found in this search.

\section{Digital and Audio Resources}
None found in this search. The search summary mentioned a UNAM Morelia course manual, \emph{Ju je uantani}, for P'urhepecha I--IV, but no URL for it appeared in the results.

\section{Curriculum Development Context}
Purepecha has a federal bilingual-education footprint, with SEP textbooks and a growing scholarly literature on a school grammar. A curriculum could build on those, after verifying the contents of the SEP documents and the dialect they use.
\end{document}